%%%%%%%%%%%%
%
% $Autor: Wings $
% $Datum: 2019-03-05 08:03:15Z $
% $Pfad: MainFunction.tex $
% $Version: 4250 $
% !TeX spellcheck = en_GB/de_DE
% !TeX encoding = utf8
% !TeX root = manual 
% !TeX TXS-program:bibliography = txs:///biber
%
%%%%%%%%%%%%

\chapter{Main Function}

Das Radarsystem dient der automatisierten Überwachung und digitalen Abbildung des unmittelbaren Umfelds. Nach der Aktivierung führt das System folgende Kernfunktionen autonom aus:

\begin{itemize}
	\item \textbf{Automatisierte Umfeldabtastung:} Durch eine koordinierte Ansteuerung des Servomotors wird der Ultraschallsensor in einem Schwenkbereich von $180^\circ$ bewegt. Diese Bewegung ermöglicht eine lückenlose Erfassung des vorderen Halbraums.
	\item \textbf{Akustische Distanzmessung:} Während des Scan-Vorgangs emittiert der Sensor hochfrequente Ultraschallimpulse. Durch die Auswertung der Echo-Laufzeit berechnet das System in Echtzeit die Distanz zu Objekten im Erfassungsbereich.
	\item \textbf{Datenübermittlung und Visualisierung:} Die ermittelten Messwerte (Winkelposition und Distanz) werden über die serielle Schnittstelle an das Host-System übertragen. Dort erfolgt die grafische Aufbereitung, welche Hindernisse ortsgetreu darstellt und ein digitales Abbild der Umgebung generiert.
\end{itemize}

Das System ist auf einen wartungsfreien Dauerbetrieb ausgelegt und passt die Abtastrate automatisch an die interne Verarbeitungsgeschwindigkeit an, um eine konsistente Datenqualität zu gewährleisten.